\subsection{Arbeitsschwerpunkte}

\begin{frame}
	{Technologie COW14\slash COW16}

	\begin{itemize}
	  \item texrex-behindthecow: 25.000 Zeilen Webpage-Cleaner
	  \item diverse Normalisierungstools (Eigenentwicklungen)
	  \item satzweise Sprachfilterung mit LangID.py + Heuristik
	  \item Annotation mit TreeTagger, FreeLing, Mate, Marmot, MaltParser, Stanford NER, Berkeley Parser\\
	    mit topologischem Modell von Cheung \& Penn (2009) usw.
	  \item Nachverarbeitungsketten standardisiert

	  \vspace{0.5cm}
	   
	  \item Und vor allem: Unmögliches ist möglich dank\\
		massiver Parallelisierung auf HPC der FU Berlin.
	\end{itemize}
\end{frame}

\begin{frame}
	{Ausführlich zur Methode}
	\centering
	\includegraphics[height=0.7\textheight]{graphics/wcc}
\end{frame}

\begin{frame}
  {COW14A-Fakten}
  \begin{center}
    \begin{tabular}[h]{lrrr}
      \hline
            & G tokens & M sentences & M documents \\
      \hline
      \hline
      DE & 20.5 & 807 & 17.1 \\
      EN & 16.8 & 608 & 9.2 \\
      ES & 7.3 & &      5.0 \\
      NL & 6.8 & 312 &  5.0 \\
      FR & 11  & 329 &  6.9 \\
      SV & 8.6 & 406 &  6.4 \\
      \hline
    \end{tabular}
  \end{center}
\end{frame}


